\chapter{Magnetic-doublet criterion (menu 16)}
\label{ch:menu16}
\menulabel{16}

\section{Purpose}

The empirical metric a Dy(III) single-molecule-magnet study reads off the
Kramers-doublet ladder: which doublet can still act as a step of the relaxation
barrier, and which one opens quantum tunnelling (QTM). For a doublet with
principal $g$ values $g_1 \le g_2 \le g_3$ and $\theta_3$ the angle (in
\emph{degrees}) between its $g_3$ axis and the ground doublet's $g_3$ axis,
\[
g_T = \frac{g_1 + g_2 + g_3 \sin\theta_3}{3},
\qquad\text{line:}\quad g_T\,\theta_3 = 20 .
\]
Below the line the doublet supports excitation above it; at or above it, QTM is
opened. The menu is pure post-processing --- the criterion needs nothing but the
doublets' $g$ tensors --- so it runs from a small JSON table.

\section{What you need}

A JSON table of the doublets: per doublet the three principal $g$ values and
\emph{either} \code{theta3} (as published tables give it) \emph{or}
\code{axis3}, the unit axis of the largest value, from which the angle to the
reference doublet's axis is computed here. The schema is in
Chapter~\ref{ch:formats}.  Alternatively, the file may be the raw text output
of an OpenMolcas \code{SINGLE\_ANISO} run (recognized by its pseudospin
banner): the $g$ tensors of every parsed multiplet are converted into the same
table, so the engine printout and the hand-written table enter the criterion
by one door.

\section{Walkthrough}

\lstinputlisting[style=fbkcode, firstline=21, lastline=38,
  title={The menu-16 example session (the source's own calibration row for
  [Dy(BC$_4$Ph$_5$)$_2$]$^-$)}]
  {examples/expected/m16_doublets.txt}

\section{What you get}

The verdict table on screen and a report file
(\file{<table>.magnetic.fbk.md}) that also carries the criteria, the
applicability domain, the known outlier and the citations.

\section{How to read the numbers}

\begin{readbox}
\begin{itemize}
  \item \textbf{The angle is in degrees} --- a measured fact of the calibration
    table, not a convention: recomputing its $g_T$ values in radians disagrees
    in the second digit. The regression check pins all 38 tabulated doublets
    down from the raw columns of the source's Table S2.
  \item \textbf{The reference doublet} (marked \code{*}) is the one whose
    $g_3$ axis is the quantisation axis the computed angles are measured
    against; its own $\theta_3$ is 0 by definition. Give the angle directly on
    every doublet and no axis is needed at all --- the numbers are then used as
    given (that is how a published table reads). An angle computed from axes is
    folded to $[0, 90]$ degrees, because a principal axis's sign is arbitrary
    (the source's own values run to $89.9^\circ$).
  \item \textbf{The criterion is empirical, not derived.}  The source states
    that the $g_1/g_2$ weights should in principle be reduced as $g_3$ rotates
    into the plane, but those historic calculations do not carry that
    information; and the line's physical background --- the material's internal
    dipolar field --- varies with packing and dilution, so the line is not a
    universal constant.
  \item \textbf{Applicability domain}: the line was calibrated on 19
    mononuclear Dy(III) complexes (2016--2025, one consistent methodology:
    9-in-7 active space, ANO-RCC three-tier basis, SA-CASSCF-SO in
    (Open)Molcas). Other ions or nuclearities would need their own calibration.
  \item \textbf{A known failure mode the criterion cannot see}: in the
    source's own set exactly one doublet is misread ---
    [Dy(Cp\textsuperscript{ttt})\textsubscript{2}]\textsuperscript{+}'s QTM
    doublet falls deep in the safe zone (product 2.52), and the source offers
    no explanation. A ``safe'' verdict is therefore the criterion's reading,
    not a proof of barrier behaviour. On the calibration table the line sits in
    the measured gap between 15.31 (largest safe product) and 30.25 (smallest
    definite-QTM product).
\end{itemize}
\end{readbox}

\section{Worked example}

The example input (\file{examples/generate\_inputs.py}) is the source's
calibration row for [Dy(BC$_4$Ph$_5$)$_2$]$^-$: the two doublets the published
table records, with $\theta_3$ as given.  The run reads the source's own
arithmetic back --- $g_T = 0.782$ and $2.993$ (the table's values), products
7.44 and 52.62 --- and places the second doublet across the line: the first
supports excitation, the second facilitates QTM.  That is the menu's purpose in
one line: given the doublet table, say which side of the line each doublet is
on, with the domain and the outlier named.

A second, engine-side example closes the same loop through a real calculation:
the fixtures' Dy(III) $4f^9$ single-ion chain
(\file{fixtures/openmolcas/dy\_smoke.out}, OpenMolcas v26.06, the input shipped
beside it) ends in a \code{SINGLE\_ANISO} block whose pseudospin multiplet is
read directly --- $g_1 = 0.913$, $g_2 = 1.991$, $g_3 = 13.144$, $\theta_3 = 0$
(a single multiplet is its own reference), $g_T = 0.968$, product $0.00$: the
ground doublet supports excitation.  The bare ion has no ligand field, so the
number carries the degeneracy of the free-ion $J = 15/2$ manifold rather than a
molecular easy axis; the fixture's role is to exercise the reading chain, and
the fixture README states that boundary.  The same output also carries the
temperature-dependence table of the magnetic susceptibility (301 points here),
and the menu quotes its 300~K entry beside the criterion line ---
$\chi T = 14.2177$~cm$^3$K/mol against the free-ion $^6$H$_{15/2}$ Curie value
of about 14.2 --- so the g-tensor reading and the bulk-susceptibility reading
of the same run are shown together.

\section{Refusals and related sections}

\begin{refusalbox}
\begin{itemize}
  \item A doublet without three finite $g$ values, an angle outside
    $[0, 180]$ degrees, a zero-length axis, a doublet with neither
    \code{theta3} nor \code{axis3} (when no reference axis is available), or a
    reference doublet that carries a nonzero \code{theta3} while providing the
    axis --- each is refused with a next step.
  \item The menu reads a table; it does not compute $g$ tensors. Those come
    from your own magnetic-property calculation --- for the OpenMolcas chain
    menu~\menu{3} generates, the doublet $g$ values are read from the
    SINGLE\_ANISO output, and this menu is the chain's post-processing step.
  \item The line is printed as source evidence with its domain and its
    outlier; the menu does not extend the criterion to other ions or
    nuclearities.
\end{itemize}
\end{refusalbox}

Related: \menu{3} (the OpenMolcas magnetic-property chain the table comes
from), \menu{10} (the crystal-field fit over the same doublet ladder), and
Chapter~\ref{ch:criteria} for the threshold table.
